\documentclass[12pt]{article}
\usepackage[T1]{fontenc}
\usepackage[utf8]{inputenc}
\usepackage[brazil]{babel}
\usepackage[margin=1in]{geometry}
\usepackage{ragged2e}
\usepackage[normalem]{ulem}
\setlength{\parindent}{0pt}
\setlength{\parskip}{0.8\baselineskip}
\begin{document}
\justifying
\noindent Verifica-se que a questão objeto da controvérsia jurídica suscitada no recurso manejado pela parte recorrente esteve afetada no representativo da controvérsia - \textbf{Tema n. }\textbf{109} da TNU - (PU no processo n. PEDILEF 0003060-22.2006.4.03.6314/ SP),  afetado em 09/10/2013 foi julgado (trânsito em julgado em 19/05/2014), por meio do qual foi elucidada a seguinte questão:

\noindent \textit{"Saber se é aplicável o art. 1º-F da Lei n. 9.494/97, com redação dada pela Lei n. 11.960/2009, em valores em atraso até a data de expedição do precatório."}

\noindent Restou firmada a seguinte tese:

\noindent \textit{"Tese firmada no Tema 810/STF: 1) O art. 1º-F da Lei nº 9.494/97, com a redação dada pela Lei nº 11.960/09, na parte em que disciplina os juros moratórios aplicáveis a condenações da Fazenda Pública, é inconstitucional ao incidir sobre débitos oriundos de relação jurídico-tributária, aos quais devem ser aplicados os mesmos juros de mora pelos quais a Fazenda Pública remunera seu crédito tributário, em respeito ao princípio constitucional da isonomia (CRFB, art. 5º, caput); quanto às condenações oriundas de relação jurídica não-tributária, a fixação dos juros moratórios segundo o índice de remuneração da caderneta de poupança é constitucional, permanecendo hígido, nesta extensão, o disposto no art. 1º-F da Lei nº 9.494/97 com a redação dada pela Lei nº 11.960/09; e 2) O art. 1º-F da Lei nº 9.494/97, com a redação dada pela Lei nº 11.960/09, na parte em que disciplina a atualização monetária das condenações impostas à Fazenda Pública segundo a remuneração oficial da caderneta de poupança, revela-se inconstitucional ao impor restrição desproporcional ao direito de propriedade (CRFB, art. 5º, XXII), uma vez que não se qualifica como medida adequada a capturar a variação de preços da economia, sendo inidônea a promover os fins a que se destina." (Tema 109 TNU).}

\noindent Com efeito, depreende-se da leitura do voto condutor do acórdão que o mesmo se encontra em conformidade com o decidido no  tema pela TNU.

\noindent A proposito, consta do acordao hostilizado: (...) Fornecimento de medicamentos, Pública, DIREITO DA SAÚDE, Fraldas, Fornecimento de insumos, Pública, DIREITO DA SAÚDE

\noindent Nessas condições, o conteúdo do Acórdão recorrido é consentâneo com o entendimento da TNU, firmado em julgamento de recurso representativo de controvérsia, o que atrai a incidência da regra prevista pelo art. 14, III, \textit{b,} da Resolução CJF n. 586/2019:

\noindent \textit{\textbf{Art. 14.}}\textit{ Decorrido o prazo para contrarrazões, os autos serão conclusos ao magistrado responsável pelo exame preliminar de admissibilidade, que deverá, de forma sucessiva: (...) }\textit{\textbf{III }}\textit{- negar seguimento a pedido de uniformização de interpretação de lei federal interposto contra acórdão que esteja em conformidade com entendimento consolidado: }\textit{\textbf{b)}}\textit{ em }\uline{\textit{\textbf{recurso representativo de controvérsia pela Turma Nacional de Uniformização}}}\textit{ ou em pedido de uniformização de interpretação de lei dirigido ao Superior Tribunal de Justiça;(...).}

\noindent Posto isso, com arrimo no \uline{\textbf{art. 14, III, }}\uline{\textit{\textbf{b}}}\textit{,} da Resolução CJF n. 586/2019 c/c art. 1.030, I, do CPC, \textbf{NEGO SEGUIMENTO }\textbf{a}o\textbf{ PU}\textbf{.}

\noindent Intimem-se. Aguarde-se o decurso do prazo recursal. Após, certifique-se o trânsito em julgado e devolva-se ao juízo de origem.
\end{document}
